\subsection{Solution to Problem 23: Simple Harmonic Motion}

\begin{solution}
The particle oscillates between $0.05$ m and $0.17$ m, so the amplitude is half the total range:
\[
A = \frac{0.17 - 0.05}{2} = 0.06 \text{ m}.
\]
The frequency is $f = 40$ Hz, so the angular frequency is
\[
\omega = 2\pi f = 80\pi \text{ rad/s}.
\]
For simple harmonic motion, the maximum acceleration occurs at the endpoints and has magnitude
\[
a_{\max} = \omega^2 A = (80\pi)^2 \times 0.06 = 6400\pi^2 \times 0.06 = 384\pi^2 \text{ m/s}^2.
\]
The maximum force on the mass is $F_{\max} = ma_{\max}$, so with $m = 0.8$ kg,
\[
F_{\max} = 0.8 \times 384\pi^2 = 307.2\pi^2 \approx 3032 \text{ N}.
\]
To the nearest newton, $F_{\max} \approx 3032$ N.
\end{solution}

\begin{takeaways}
\begin{itemize}
    \item SHM amplitude: $A = \dfrac{\text{max displacement} - \text{min displacement}}{2}$
    \item Maximum acceleration: $a_{\max} = \omega^2 A$ where $\omega = 2\pi f$
    \item Newton's second law: $F_{\max} = ma_{\max}$ at maximum acceleration
\end{itemize}
\end{takeaways}
